\newcounter{interludecounter}
\newcommand{\BookSpecialKicker}{}
\newcommand{\BookTitleSep}{\texorpdfstring{\quad}{\space}}

\newcommand{\interlude}[1]{%
  \stepcounter{interludecounter}%
  \renewcommand{\BookSpecialKicker}{INTERLUDE~\Roman{interludecounter}}%
  \chapter*{间幕 \Chinese{interludecounter}\BookTitleSep #1}%
  \addcontentsline{toc}{chapter}{间幕 \Chinese{interludecounter}\BookTitleSep #1}%
  \markboth{间幕 \Chinese{interludecounter}\quad #1}{间幕 \Chinese{interludecounter}\quad #1}%
  \renewcommand{\BookSpecialKicker}{}%
}

\newcommand{\introduction}[1]{%
  \renewcommand{\BookSpecialKicker}{INTRODUCTION}%
  \chapter*{导言\BookTitleSep #1}%
  \addcontentsline{toc}{chapter}{导言\BookTitleSep #1}%
  \markboth{导言\quad #1}{导言\quad #1}%
  \renewcommand{\BookSpecialKicker}{}%
}

\newcommand{\bookbackchapter}[2]{%
  \renewcommand{\BookSpecialKicker}{#1}%
  \chapter*{#2}%
  \addcontentsline{toc}{chapter}{#2}%
  \markboth{#2}{#2}%
  \renewcommand{\BookSpecialKicker}{}%
}

\newcommand{\chaptersubtitle}[1]{%
  \begin{flushleft}
    \vspace{-10pt}%
    {\sffamily\fontsize{10pt}{14pt}\selectfont\color{BookGray}#1\par}%
    \vspace{12pt}%
  \end{flushleft}%
}

\newenvironment{chapteropeningnote}{%
  \begin{tcolorbox}[
    enhanced,
    breakable,
    colback=BookPale,
    colframe=MadokaPink,
    boxrule=0pt,
    borderline west={1.2pt}{0pt}{MadokaPink},
    left=7pt,right=7pt,top=6pt,bottom=6pt,
    before skip=4pt,after skip=14pt,
    fontupper=\fontsize{\BookOpeningNoteFontSize}{\BookOpeningNoteLineHeight}\selectfont,
    title={题记／译者案},
    coltitle=white,
    fonttitle=\sffamily\bfseries
      \fontsize{\BookOpeningNoteLabelFontSize}{\BookOpeningNoteLabelLineHeight}\selectfont
  ]
}{\end{tcolorbox}}

% Backward-compatible wrapper; new long notes should use chapteropeningnote.
\newcommand{\chapternote}[1]{%
  \begin{chapteropeningnote}#1\end{chapteropeningnote}%
}

\newcommand{\translatornote}[1]{\footnote{\textbf{译注：}#1}}

\newcommand{\bookpart}[1]{%
  \cleardoublepage
  \refstepcounter{part}%
  \phantomsection
  \addcontentsline{toc}{part}{第\Chinese{part}部分\BookTitleSep #1}%
  \markboth{第\Chinese{part}部分\quad #1}{}%
  \thispagestyle{empty}%
  \begin{tikzpicture}[remember picture,overlay]
    % Part pages are always recto: shift to the bound text-block centre.
    \node[opacity=.94] at
      ($(current page.center)+(\BookRectoTextCenterShift,17mm)$) {\BookPartMotif};
    \draw[MadokaPink,line width=.5pt]
      ($(current page.south west)+(14mm,20mm)$)--($(current page.south east)+(-14mm,20mm)$);
  \end{tikzpicture}%
  \vspace*{69mm}%
  \begin{center}
    {\sffamily\bfseries\fontsize{9pt}{11pt}\selectfont\color{BookGray}%
      PART~\Roman{part}\par}%
    \vspace{5pt}%
    {\sffamily\bfseries\fontsize{22pt}{28pt}\selectfont #1\par}%
    \vspace{8pt}%
    {\color{HomuraViolet}\runeecho{PART \Roman{part}}\par}%
  \end{center}%
  \clearpage
  \thispagestyle{empty}\null\clearpage
}

% Future plate contract: exactly four single-page 131 x 182 mm bleed PDFs.
\newcommand{\BookRequireSinglePagePlate}[1]{%
  \IfFileExists{#1}{%
    \ifnum\XeTeXpdfpagecount "#1"=1\relax\else
      \PackageError{the-very-soil}{Plate #1 is not a single-page PDF}%
        {Split the artwork into page-01.pdf through page-04.pdf.}%
    \fi
  }{%
    \PackageError{the-very-soil}{Missing plate #1}%
      {Each enabled Part plate unit must contain exactly four page PDFs.}%
  }%
}

\newcommand{\includepartplates}[1]{%
  \foreach \n in {01,02,03,04}{%
    \BookRequireSinglePagePlate{assets/plates/#1/page-\n.pdf}%
  }%
  \foreach \n in {01,02,03,04}{%
    \includepdf[
      pages=1,
      noautoscale=true,
      trim=3mm 3mm 3mm 3mm,
      clip,
      pagecommand={\thispagestyle{empty}}
    ]{assets/plates/#1/page-\n.pdf}%
  }%
}
